\subsection{A cubic pencil refuting the literal middle fixed-part clause}
\paragraph{J6C (Refuted by an explicit cubic pencil).}
At 44:4, Theorem~6 says, under the conditions of Theorem~5,
``si $2i=n$ le facteur fixe de $P^i(Y)$ est identique à
l'image de $P^i(X)$''. The inserted condition on $Y$ in
Theorem~5 is ``section hyperplane générique d'un pinceau''
whose fibres are geometrically integral. No smoothness of the
pencil axis is imposed in either supplied transcription.
The word ``générique'' modifies the hyperplane section of the
pencil; the pencil itself is not required to be general.
The following example meets those literal hypotheses and has
zero source but a positive-dimensional fixed target factor.
The cycle hypotheses in Theorem~5 can, if needed, be met by
taking $Z=Z'=0$ with $i=j=2$.

\emph{Existing-source check.}
The underlying elliptic-cone geometry is classical:
\cite[Lemma 8.1, pp.~314--315]{DGJ_ClemensGriffiths1972}
describes the cone of lines through an Eckardt point, and
\cite[Proposition 10, p.~384; Lemma 29, p.~398]{DGJ_Roulleau2009}
studies the associated elliptic curve and homomorphisms on the
complex Fano surface. The explicit pencil and positive-characteristic
calculation below are supplied with a complete proof; no priority
claim is made for this construction.

Let $k=\overline{\mathbb F}_5$, and let $X\subset\mathbb P^5_k$
be the cubic fourfold
\begin{equation}\label{eq:J6C-fourfold}
 x_0^3+x_1^3+x_2^3+x_3^2x_4+x_4^3+x_5^3=0.
\end{equation}
Consider the hyperplane pencil $x_5=t x_4$, with parameter
$t\in\mathbb A^1_k$, and the member $x_4=0$ at infinity.
Its axis is the integral surface
\[
 B:\quad x_4=x_5=0,\qquad x_0^3+x_1^3+x_2^3=0.
\]
It is the cone over the smooth plane cubic
\[
 E:\quad u_0^3+u_1^3+u_2^3=0\quad\text{in }\mathbb P^2_k,
\]
with vertex $p=[0:0:0:1:0:0]$. Choose
$e_0=[1:-1:0]$ on $E$ as origin. Let $K=k(t)$, and let
$Y/K$ be the generic pencil member. The bundle
$\mathcal O_X(1)$ is the very ample bundle of this displayed
embedding. The morphisms below are all defined over $K$.

\emph{Smooth total variety and integral pencil members.}
The gradient of \eqref{eq:J6C-fourfold} is
\[
 (3x_0^2,3x_1^2,3x_2^2,2x_3x_4,
   x_3^2+3x_4^2,3x_5^2).
\]
A common zero forces $x_0=x_1=x_2=x_5=0$, then
$x_3x_4=0$ and $x_3^2+3x_4^2=0$, and finally
$x_3=x_4=0$. Thus $X$ is smooth. For a finite parameter,
the member is the cubic in $\mathbb P^4$ with equation
\begin{equation}\label{eq:J6C-member}
 x_0^3+x_1^3+x_2^3+x_3^2x_4+a x_4^3=0,
 \qquad a=1+t^3.
\end{equation}
If $a\ne0$, the same derivative calculation shows that it
is smooth. If $a=0$, its singular locus has the single
geometric point $[0:0:0:0:1]$. The member at infinity is
the cone, with vertex $p$, over the smooth Fermat cubic
surface $x_0^3+x_1^3+x_2^3+x_5^3=0$; its singular locus
also has just one point.

Every member is geometrically integral. Indeed, over an
algebraically closed field a reducible hypersurface in
$\mathbb P^4$ has two components meeting in dimension at
least two, and their intersection is singular for the
product equation. A repeated component is singular along
a three-dimensional subset. Both possibilities contradict
the computed empty or zero-dimensional singular locus.
This also proves geometric integrality of the generic
member after extending $K$ to an algebraic closure. The
three finite singular parameters are distinct, since
$1+t^3$ is relatively prime to $3t^2$ in characteristic five.

The pencil is the incidence family in $X\times\mathbb P^1$.
Because $x_4,x_5$ form a regular sequence along the
codimension-two axis, the two-generated ideal is of linear
type \cite[Proposition IV-25, pp.~172--173]{DGJ_EisenbudHarris2000};
its blowup is that incidence family. Thus its scheme
fibres are exactly the displayed hyperplane sections,
with no additional fibre components. Equivalently, each
local incidence equation restricts to a nonzerodivisor on
each fibre of $X\times\mathbb P^1$, so this is a flat
hyperplane family. The singular axis is the sole missing
hypothesis relative to the positive theorem J6P.

\emph{The source representative vanishes.}
Weak Lefschetz for the smooth ample hypersurface
$X\subset\mathbb P^5$ gives
\[
 H^3(X,\mathbb Q_\ell)=H^3(\mathbb P^5,\mathbb Q_\ell)=0
 \qquad(\ell\ne5).
\]
Consequently
\begin{equation}\label{eq:J6C-zero-source}
 \operatorname{Ab}^2(X)=0,\qquad P^2_{\rm inc}(X)=0.
\end{equation}
For the first conclusion use the miniversal-cycle injection
\cite[Proposition 8.2, p.~31]{DGJ_ACMVFunctor}; for the
second use Saito's incidence bound
\cite[Proposition 3.4, pp.~109--110]{DGJ_Saito1979}.
Weak Lefschetz is cited in
\cite[(4.1.6), pp.~218--219]{DGJ_Deligne1980}.

\emph{A constant elliptic factor in the target.}
Blow up the $K$-point $p$ of the smooth threefold $Y$:
\[
 \pi:\widehat Y=\operatorname{Bl}_pY\longrightarrow Y,
 \qquad D\simeq\mathbb P^2_K.
\]
Let $i:S\hookrightarrow\widehat Y$ be the strict transform
of $B_K$. This is the smooth ruled resolution of the
elliptic cone, with projection $q:S\to E_K$. Concretely,
its fibre over $e=[u_0:u_1:u_2]$ is the strict transform
of the ruling line
\[
 [\lambda u_0:\lambda u_1:\lambda u_2:\mu:0:0]
 \quad\text{in }B_K.
\]
Thus $S$ is a $\mathbb P^1$-bundle over $E_K$, and every
ruling fibre $\ell$ satisfies
$\pi^*H\cdot\ell=1$ and $D\cdot\ell=1$.

The multiplicity of $B_K$ as a divisor of $Y$ at $p$ is
three. To check it, use $x_3=1$ and put $v=x_4$; equation
\eqref{eq:J6C-member} becomes
\[
 v+a v^3+u_0^3+u_1^3+u_2^3=0.
\]
The three $u_j$ are regular parameters at $p$, and
$v(1+a v^2)=-(u_0^3+u_1^3+u_2^3)$, with the first factor
in parentheses a unit. Hence $v$ has order three. As a
Cartier divisor on $\widehat Y$, therefore,
\[
 S=\pi^*H-3D,
 \qquad
 \deg(N_{S/\widehat Y}|_\ell)=1-3=-2.
\]

Put $J=\operatorname{Pic}^0(E)$, a constant elliptic
curve over $k$. The algebraic-correspondence operations
\[
 i_*q^*:A^1(E_K)\longrightarrow A^2(\widehat Y),\qquad
 q_*i^*:A^2(\widehat Y)\longrightarrow A^1(E_K)
\]
preserve algebraic equivalence and regularity. Here $i^*$
is the refined pullback for the Cartier embedding $i$;
$q_*$ sends zero-cycles on the surface $S$ to divisor
classes on $E_K$. They induce homomorphisms of abelian
varieties
\[
 \widehat u:J_K\longrightarrow\operatorname{Ab}^2(\widehat Y),
 \qquad
 \widehat v:\operatorname{Ab}^2(\widehat Y)\longrightarrow J_K.
\]
The Cartier self-intersection formula
\cite[Definition 42.29.1(1)(b)]{DGJ_StacksGysin}
and the projection formula
\cite[Lemma 42.26.4]{DGJ_StacksProjection} give,
for every $L\in\operatorname{Pic}^0(E_K)$,
\begin{align*}
 q_*i^*i_*q^*L
 &=q_*\bigl(c_1(N_{S/\widehat Y})\cdot q^*L\bigr)\\
 &=\bigl(q_*c_1(N_{S/\widehat Y})\bigr)\cdot L=-2L.
\end{align*}
Thus
\begin{equation}\label{eq:J6C-minus-two}
 \widehat v\widehat u=[-2]_{J_K}.
\end{equation}
This sign is an equality of divisor-class homomorphisms
on $J$, not an unstated choice of sign for an elliptic
principal polarization.

The point-blowup formula gives
\[
 \operatorname{CH}^2(\widehat Y)
    =\pi^*\operatorname{CH}^2(Y)\oplus\mathbb Z[\text{line in }D].
\]
The exceptional integer is constant under algebraic
equivalence, so $\pi^*$ and $\pi_*$ are inverse on $A^2$.
They also preserve regular families and therefore induce
inverse isomorphisms of the universal regular
representatives. This is the $r=3$ case of the blowup
formula \cite[Theorem 2.8(3)]{DGJ_Jiang2023}.
Define
\[
 u=\pi_*\widehat u:J_K\longrightarrow\operatorname{Ab}^2(Y),
 \qquad
 v=\widehat v\pi^*:\operatorname{Ab}^2(Y)\longrightarrow J_K.
\]
Equations above give $vu=[-2]_{J_K}$. In characteristic
five this map is an isogeny of degree four. Hence $u$ has
finite kernel and a one-dimensional abelian image. Since
its source is constant and $u$ is defined over $K$, the
universal property of the $K/k$-trace puts this image in
the fixed part of $\operatorname{Ab}^2(Y)$. That fixed
part is therefore positive-dimensional, contradicting
\eqref{eq:J6C-zero-source}.

The factor is detected by incidence equivalence as well:
$q_*i^*\pi^*$ is an actual correspondence from $Y$ to the
curve $E_K$, taking codimension-two classes to divisor
classes. It is one of the Picard tests defining the
incidence quotient. Thus the constant family $i_*q^*$
cannot vanish in that quotient, since composing with
this test is $[-2]$. In particular the counterexample
does not depend on the disputed integral autoduality
of a universal representative.

\emph{Conclusion and remaining scope.}
For $n=4,i=2$, a generic hyperplane section of this
integral-fibre pencil has a nonzero fixed elliptic
factor, although the ambient fourfold has zero
codimension-two representative. The literal middle
clause of Theorem~6 is therefore false. This does not
refute J6P, whose smooth-axis hypothesis is essential,
or settle the separate stable- and high-range assertions
in all codimensions. The exact gradient ideals,
singular supports, origin, parameter roots, and normal
degree are checked by \texttt{check\_cubic\_trace.py};
the constant-factor conclusion is the correspondence
proof above, not an output of the finite computation.
